% Section 6 -- Evaluation

\subsection{Experimental Setup}

We evaluate on Twitter (41.6\,M nodes, 1.46\,B edges), 16 partitions,
buffer capacity 4, and two \SI{24}{GB} RTX A5000 GPUs connected over
PCIe. We use the existing GE\textsuperscript{2} Twitter
DistMult/Adagrad configuration and compare three runs:
\textsc{Lane\_Matched} with peer execution disabled,
\textsc{Lane\_Matched} with peer execution enabled, and
\textsc{Disjoint\_Rounds} with peer execution enabled.

For the peer-enabled runs we set
\texttt{GEGE\_STATEFLOW\_PEER\_RUNTIME\_SCOPE=embeddings}, so the
runtime executes peer transfers for node embeddings while optimizer
state remains on the existing host path. Under this scope, executed
peer bytes should equal exactly half of the compiler-reported planned
peer bytes. We report compile-time selection cost, admits, and handoff
distribution from the IR, then runtime epoch time, executed peer
bytes, host fallback bytes, descriptor mismatches, and peer-sync wait
time from the runtime telemetry.

\subsection{Compile-Time Effect}

Table~\ref{tab:ir:example} summarizes the two compiled plans.
\textsc{Lane\_Matched} lowers selection cost from
$1.499 \times 10^{11}$ to $1.291 \times 10^{11}$ (14\%), reduces host
admits from 72 to 62, and changes the handoff distribution from
\{\texttt{FULL\_RELOAD}:10, \texttt{ROTATING\_OVERWRITE}:8,
\texttt{PEER\_RELAY}:8\} to
\{\texttt{FULL\_RELOAD}:0, \texttt{ROTATING\_OVERWRITE}:18,
\texttt{PEER\_RELAY}:18\}. The important point is that this entire
improvement is visible before runtime: the compiler eliminates every
cold admit from the plan and exposes \SI{37.49}{GB} of peer-relay
opportunity, versus \SI{16.66}{GB} for \textsc{Disjoint\_Rounds}.

\subsection{Runtime Effect}

\begin{table}[t]
\centering
\small
\begin{tabular}{lrrrrr}
\toprule
Variant & Epoch (s) & Peer GB & Copies & Fallback GB & Mismatch \\
\midrule
\textsc{Lane\_Matched}, peer off          & 208.256 & 0.00  & 0  & 0.00 & 0 \\
\textsc{Lane\_Matched}, peer on (emb.)    & 209.450 & 18.74 & 18 & 0.00 & 0 \\
\textsc{Disjoint\_Rounds}, peer on (emb.) & 209.525 & 8.33  & 8  & 0.00 & 0 \\
\bottomrule
\end{tabular}
\caption{Twitter 16-partition runtime on two RTX A5000s. Peer-on runs
use \texttt{GEGE\_STATEFLOW\_PEER\_RUNTIME\_SCOPE=embeddings}, so the
executed peer bytes are exactly half of the planned peer bytes in
Table~\ref{tab:ir:example}.}
\label{tab:eval:runtime}
\end{table}

Table~\ref{tab:eval:runtime} gives the runtime result. For
\textsc{Lane\_Matched}, the peer-enabled run executes all 18 planned
embedding handoffs, totaling \SI{18.74}{GB} of device-to-device
traffic, with zero host fallback and zero descriptor mismatch. The
runtime telemetry reports \texttt{peer\_sync\_wait\_ms=1019.646} for
the epoch. For \textsc{Disjoint\_Rounds}, the runtime executes all 8
planned embedding handoffs, totaling \SI{8.33}{GB}, also with zero
fallback and zero mismatch.

\textsc{Lane\_Matched} remains the best peer-enabled schedule, but only
by \SI{0.075}{s}: \SI{209.450}{s} versus \SI{209.525}{s}. Both peer-on
runs are about \SI{1.2}{s} slower than the
\textsc{Lane\_Matched}/peer-off baseline at \SI{208.256}{s}. On this
PCIe A5000 setup, peer execution therefore displaces host traffic
correctly but does not convert that displacement into an epoch-time
speedup. The implication is that the remaining ceiling is the swap
runtime itself rather than descriptor quality: the compiler removes
cold reloads and doubles executed peer traffic, but the wall-clock path
is still dominated by swap orchestration overhead on this hardware.

\subsection{Takeaway}

The runtime result is still valuable. Phase 4b is
\emph{correctness-complete}: the compiler-emitted descriptors survive
projection, validation, slot alignment, and swap-time execution with
zero fallback. The peer-enabled runtime also preserves the compiler's
ordering preference: \textsc{Lane\_Matched} executes more than twice the
peer traffic of \textsc{Disjoint\_Rounds}
(\SI{18.74}{GB} vs.\ \SI{8.33}{GB}) and remains the best peer-enabled
variant. What the result does \emph{not} show is a wall-clock win on
this hardware. The next step is therefore Phase 5 overlap and
staggering, or evaluation on a faster peer interconnect.
